%==============================================================================
\section{Exact methods}\label{sec:exact}
%==============================================================================

Section~\ref{sec:gap} showed that the coverage bound $q=\lvert\U\rvert-b$ equals the
optimum on a large part of the instance space. Section~\ref{sec:sota} showed that the
strongest of the compact formulations reimplemented here has a linear relaxation equal to that
same quantity. This section asks whether an exact method can do better, and builds
two that are designed to.

\paragraph{Why these two methods.}
A compact integer formulation is the natural first recourse, and three from the
literature serve here as baselines. Their relaxations do not exceed the coverage
bound, so a branch-and-cut built on any of them is fast where that bound is tight and
adrift where it is not. Escaping this requires a method that does not rest on a
compact relaxation, and two classical routes qualify, for opposite reasons.

Benders decomposition is the textbook fit when the subproblem is polynomially
solvable \citep{Rahmaniani2017}. The loading problem is such a subproblem, so it
becomes an oracle returning exact cuts, and the master's bound is assembled from true
subproblem values rather than from a linear relaxation of the loading. The algorithm
never forms that relaxation at all --- which is a statement about the loading and not
about the master, whose own root bound Observation~\ref{obs:rootq} finds sitting on the
coverage row on $87.2\%$ of recorded runs. This places the solver in the lineage of logic-based
\citep{Hooker2003} and combinatorial \citep{Codato2006} Benders decomposition.

Branch-and-price attacks the same weakness from the other side. Reformulating over
configurations, its position--transition variant carries a relaxation that can provably
exceed the coverage bound (Proposition~\ref{prop:ptf}), which is the one construction in
this report that does. How often it does so on real instances is a separate question,
and Section~\ref{ssec:results-bnp} answers it.

\begin{table}[t]\centering
\footnotesize
\begin{tabular}{@{}llll@{}}
\toprule
Method & Source & Formulation & Linear relaxation \\
\midrule
F4          & \citet{Catanzaro2015} & compact, arc-based            & in general below $q$ \\
LSS         & \citet{Laporte2004}   & compact, tool-state           & in general below $q$ \\
SSPMF       & \citet{daSilva2024}   & multicommodity flow           & $=q$ when every tool is used \\
BBC         & this work             & branch-and-Benders-cut        & no relaxation of the loading \\
PCF, PCF$'$ & this work             & position--configuration B\&P  & $=q$ with the counting rows \\
PTF         & this work             & position--transition B\&P     & $\ge q$, and can exceed it \\
\bottomrule
\end{tabular}
\caption{The three compact baselines and the methods of this report, ordered by the
strength of their linear relaxation relative to the coverage bound $q$. Every method
is exact; they differ in whether, and by what means, they get past that bound. The
entries for F4 and LSS are hedged because the source states only that earlier models
fall below $q$ in the great majority of cases \citep{daSilva2024}, not in all.}
\label{tab:methods}
\end{table}

\subsection{Compact formulations and the limit they reach}\label{ssec:compact}

Three compact formulations from the literature serve as baselines: Formulation~4 of
\citet{Catanzaro2015}, the tool-state model of \citet{Laporte2004}, and the
multicommodity-flow model of \citet{daSilva2024}. Section~\ref{ssec:sota-compact}
describes all three and records the omission of the stronger Formulation~5 as a
limitation. All three were reimplemented for this study, and
Section~\ref{ssec:validation} describes how they were checked. The last is singled out
here by its relaxation.

\begin{proposition}[\citealp{daSilva2024}]\label{prop:sspmf-lp}
The linear relaxation of the multicommodity-flow formulation equals
$\lvert T\rvert-b$, the number of tools minus the magazine capacity. Under the
standing assumption that every tool is required by some job, so that $\U=T$, this is
the coverage bound $q$ of Proposition~\ref{prop:lb2}.
\end{proposition}

The hypothesis matters and is stated explicitly by the authors. On an instance with
unused tools the two quantities differ, and it is $q$, not $\lvert T\rvert-b$, that
bounds the optimum. The instances of the computational study all satisfy the
hypothesis.

This result matters twice. It says that the strongest of the compact relaxations
considered here is no stronger than an elementary counting bound. And since $q$ equals $\Zst$ on a large
part of the instance space (Section~\ref{sec:gap}), it predicts that such a model will
be fast exactly where the bound is tight and weak where it is not.
Section~\ref{ssec:results-tight} tests that prediction directly.

\subsection{A branch-and-Benders-cut algorithm}\label{ssec:bbc}

\paragraph{The idea.}
A schedule is an order together with a loading. Once the order is fixed the minimum
loading cost is computed exactly and in polynomial time. The order therefore goes into
an integer master problem, and the loading cost is represented by a single continuous
variable $\theta$, refined by cuts generated from the exactly solved loading
subproblem. In contrast to column generation, whose pricing subproblem for this
problem is NP-hard (Section~\ref{ssec:pos}), the Benders subproblem here is
polynomial, so each cut is a valid optimality cut built from an exactly solved subproblem.

Figure~\ref{fig:benders} shows the loop.

\begin{figure}[t]\centering
\begin{tikzpicture}[font=\footnotesize,>=Stealth,
   blk/.style={draw,semithick,rounded corners,align=center,minimum height=1.6cm,text width=3.7cm}]
  \node[blk,draw=cBlue!70,fill=white] (m) {\textbf{Master} (hard)\\[2pt] choose the job order,\\ minimising the stand-in $\theta$};
  \node[blk,right=3.4cm of m,draw=cGreen!70,fill=white] (s) {\textbf{Subproblem} (easy)\\[2pt] load the tools for that\\ order, \emph{exactly}};
  \draw[-{Stealth[length=3mm]},semithick,black!60] (m.north east) to[bend left=20] node[above,font=\scriptsize]{an order} (s.north west);
  \draw[-{Stealth[length=3mm]},semithick,black!60] (s.south west) to[bend left=20]
      node[below,align=center,font=\scriptsize]{Benders cut:\\ $\theta\ge$ true cost} (m.south east);
  \node[below=6mm of $(m)!0.5!(s)$,text width=8.6cm,align=center,font=\scriptsize,text=black!60]
     {repeat until $\theta$ equals the loading cost; the order is then a certified optimum};
\end{tikzpicture}
\caption{The branch-and-Benders-cut loop. The master picks an order using an
optimistic stand-in $\theta$ for the loading cost. The subproblem returns the exact
loading cost of that order. If $\theta$ undershot, one cut lifts it, and that cut is
valid for every order, not only the current one. Because the subproblem is polynomial
each cut comes from an exactly solved subproblem, so the loop halts at a proven optimum. The cuts are separated
inside a branch-and-bound tree over the order variables, which is what makes the
method a \emph{branch}-and-Benders-cut.}
\label{fig:benders}
\end{figure}

\paragraph{The master problem.}
Introduce a depot node $d$ with an empty tool requirement, and arc variables
$x_{ij}\in\rset{0,1}$ for $i,j\in J\cup\rset{d}$ with $i\ne j$, where $x_{ij}=1$ means
job $j$ runs immediately after $i$. The depot marks the start and the end, so the
$x$ variables describe a Hamiltonian path through the jobs. The assignment
constraints
\[
  \sum_{i\ne j}x_{ij}=1\ \ (\forall j),\qquad \sum_{j\ne i}x_{ij}=1\ \ (\forall i),
\]
together with subtour-elimination constraints
$\sum_{i,j\in S}x_{ij}\le\lvert S\rvert-1$ for $S\subseteq J$, enforce that. The
master problem is then
\[
  \min\ \theta
\]
subject to the assignment constraints, the subtour constraints, the bound rows
described next, and the Benders cuts generated during the search.

Three families of rows bound $\theta$ from below before any cut is generated.

\emph{The pairwise bound.} At the boundary between consecutive jobs $i$ and $j$ at
least $w_{ij}:=\max(0,\lvert T_i\cup T_j\rvert-b)$ tools are switched, since the
magazine holds $b$ tools but must serve $T_i$ and then $T_j$. This is the bound of
\citet{Tang1988}. Because the subproblem charges insertions from an empty magazine,
the first job also costs at least $\lvert T_{j}\rvert$ insertions, and that term is
carried on the depot arcs. Summing gives the valid inequality
\[
  \theta\ \ge\ \sum_{i,j\in J} w_{ij}\,x_{ij}\;+\;\sum_{j\in J}\lvert T_j\rvert\,x_{d,j},
\]
in which the two contributions are disjoint by position and so may be added.

\emph{The coverage row.} Every tool of $\U$ must be loaded at least once. In the
empty-start cost this gives $\theta\ge\lvert\U\rvert$ directly. This row was absent
from the first implementation, and its absence was the reason the solver performed
poorly in the first campaign; Section~\ref{ssec:validation} records the correction.

\emph{Triplet bounds.} The pairwise argument generalises to ordered triples of jobs,
giving $O(n^{3})$ valid inequalities of the form
$\theta\ge w_{ijk}(x_{ij}+x_{jk}-1)$ with
$w_{ijk}=\max(0,\lvert T_i\cup T_j\cup T_k\rvert-b)$. These are optional and their
effect is measured in Section~\ref{ssec:results-bbc}.

Subtour separation is exact by construction: violated constraints are detected at
integer candidate solutions by a connected-component search, so no optimal solution is
ever removed. Separating them at fractional points would be an acceleration, not a
correctness requirement, and is deliberately absent from the measured configuration.

\paragraph{The subproblem.}
The subproblem must be written in the master's own variables, or the cut it produces
cannot be added to the master. Fix a candidate $\bar{x}$ and introduce, for each job
$j$ and tool $t$, a variable $y_{j,t}\in[0,1]$ for ``tool $t$ occupies the magazine
while job $j$ runs'' and $z_{j,t}\ge0$ counting an insertion of $t$ at $j$. The loading
problem is
\begin{alignat}{3}
  \min\quad & \textstyle\sum_{j}\sum_{t} z_{j,t} \notag\\
  \text{s.t.}\quad
   & y_{j,t}\ \ge\ 1
     &\qquad& t\in T_j &\qquad& [\,\nu_{j,t}\,] \label{eq:sp-req}\\
   & \textstyle\sum_{t} y_{j,t}\ \le\ b
     && \forall j && [\,\mu_j\ge0\,] \label{eq:sp-cap}\\
   & z_{j,t}\ \ge\ y_{j,t}-y_{i,t}-(1-x_{ij})
     && i\ne j,\ \forall t && [\,\lambda_{ijt}\ge0\,] \label{eq:sp-link}\\
   & z_{j,t}\ \ge\ y_{j,t}-(1-x_{dj})
     && \forall j,\ \forall t && [\,\lambda^{d}_{jt}\ge0\,] \label{eq:sp-depot}
\end{alignat}
with $y,z\ge0$. Row~\eqref{eq:sp-link} is the only place the master's variables enter:
when the master selects the arc $(i,j)$, a tool present at $j$ but absent at $i$ must be
inserted at $j$; when it does not, the row is slack. Row~\eqref{eq:sp-depot} does the
same for the first job, charging its whole requirement because the magazine starts
empty. For a fixed $\bar x$ the master's variables appear only on the right-hand side,
so what remains is a linear program in $y$ and $z$ alone; its constraint matrix has the
consecutive-ones property along the path $\bar x$ selects and is therefore totally
unimodular, so the relaxation is integral and, by Proposition~\ref{prop:ktns}, its
optimum equals $\KTNS$ of the order.

\paragraph{The cut, in full.}
Only $\bar x$ appears on the right-hand side of \eqref{eq:sp-link} and
\eqref{eq:sp-depot}, so in the dual it appears only in the objective, and the dual
feasible region does not depend on $\bar x$ at all. Writing the dual explicitly, the
constraint dual to $y_{j,t}$ is
\[
  -\mu_j-\lambda^{d}_{jt}-\sum_{i\ne j}\lambda_{ijt}+\sum_{k\ne j}\lambda_{jkt}
  +\nu_{j,t}\,[\,t\in T_j\,]\ \le\ 0 ,
\]
and the constraint dual to $z_{j,t}$ is
\[
  \lambda^{d}_{jt}+\sum_{i\ne j}\lambda_{ijt}+\eta_{j,t}\,[\,t\notin T_j\,]\ \le\ 1 ,
\]
where $\eta$ prices the non-negativity of $z$ on tools the job does not require and
carries objective coefficient zero, so it never reaches the cut. For any dual-feasible
$(\bar\lambda,\bar\lambda^{d},\bar\mu,\bar\nu)$, weak duality gives the Benders
optimality cut
\begin{equation}\label{eq:benders-cut}
  \theta\ \ge\
  \underbrace{\sum_{i\ne j}\sum_{t}\bar\lambda_{ijt}\,(x_{ij}-1)}_{\text{job-to-job arcs}}
  \;+\;\underbrace{\sum_{j}\sum_{t}\bar\lambda^{d}_{jt}\,(x_{dj}-1)}_{\text{depot arcs}}
  \;-\;b\sum_{j}\bar\mu_j
  \;+\;\sum_{j}\sum_{t\in T_j}\bar\nu_{j,t}.
\end{equation}
Collecting terms, the coefficient of $x_{ij}$ is $\sum_t\bar\lambda_{ijt}$ and every
remaining term is constant, so \eqref{eq:benders-cut} is an affine inequality in the
master's variables and can be added to it verbatim. This is the form the implementation
builds, and it is the form Section~\ref{ssec:disc-locality} analyses when it observes
that every cut the master ever receives is supported on arcs.

Two properties follow immediately. Each cut is a valid lower bound on the loading cost
of \emph{every} order, not only of $\bar x$, because the dual point is feasible for all
of them; it is tight at $\bar x$ when the dual point is optimal there. And because
\eqref{eq:benders-cut} is affine in $x$, its value at a \emph{fractional} $\bar x$ is
defined by the same expression: separating at a fractional master solution means solving
the dual with the fractional $\bar x$ in its objective and adding the resulting affine
row. No reinterpretation of the inequality is required, which is why the fractional
variant of Section~\ref{ssec:accel} needs no separate validity argument.

\begin{remark}[A term whose omission cut off optima]\label{rem:depotterm}
The depot sum in \eqref{eq:benders-cut} is non-positive, so dropping it makes the cut
\emph{tighter} rather than looser, and an implementation that omits it will look
correct on most instances. It was omitted from the first implementation. Because
$\bar\lambda^{d}$ and $\bar\nu$ trade off at zero reduced cost, degenerate dual optima
put positive weight on depot arcs that the order does not use, and the over-tight cut
then removes genuinely optimal $(x,\theta)$ points: $193$ violations were found on
random instances before the term was restored. The campaign reported in
Section~\ref{sec:comp} was run after the correction.
\end{remark}

\subsection{Strengthening the Benders solver}\label{ssec:accel}

Four strengthenings were implemented, each as an independent switch, and each checked
against brute force on small instances before any campaign run so that a strengthening
can never silently return a wrong optimum. All four are classical elements of the
toolbox surveyed by \citet{Rahmaniani2017}.

\paragraph{Fractional Benders cuts.}
The optimality cut is derived at an integer order, but nothing prevents separating the
same cut at the fractional solutions produced at a node before branching. The
intention is to attack tailing-off: an integer-only scheme learns about the loading
cost only once an order is fully committed, whereas a fractional cut extracts the same
dual information from a partial order and removes it from the relaxation at once.

Whether that pays is an empirical question, and Section~\ref{ssec:results-bbc} answers
it.

\paragraph{A strong incumbent at the root.}
A branch-and-bound search prunes in proportion to the quality of its incumbent. A
compact form of the hybrid genetic search of \citet{Mecler2021} is run once, at the
root: a population of constructive seeds, being nearest-neighbour orders on the
pairwise weights $w_{ij}$ and a largest-tool-first order, each improved by a
variable-neighbourhood local search using adjacent swaps and short segment
relocations, with every candidate scored by the exact loading oracle, evolved under
order crossover with a diversity guard.

The resulting order is used twice. As a warm start it gives the solver an incumbent
from node zero. And it seeds a Benders cut: the loading dual is solved once for the
heuristic order and its cut added before the first relaxation, which is the
initial-cut-pool acceleration of \citet{Rahmaniani2017}.

\paragraph{A conflict-graph bound.}
The master's root bound is pairwise plus coverage, and both are visible to the linear
relaxation. The combinatorial structure they miss lives in the conflict graph $G$ of
Section~\ref{ssec:clutter}, whose vertices are the jobs and whose edges join pairs
that cannot share a configuration. A clique of $G$ is a set of pairwise incompatible
jobs, which must occupy that many distinct configurations, so $\Kst\ge\omega(G)$; more
generally $\Kst\ge\chi(G)$, and where $G$ is non-bipartite one has $\chi(G)\ge3$ even
when $\omega(G)=2$. The resulting constant row
$\theta\ge(\chi_{\mathrm{lb}}-1)+\min(b,\lvert\U\rvert)$ is added whenever it strictly
exceeds the coverage row. This is the simplest of the conflict-graph inequalities studied by
\citet{Atamturk2000}. Its reach is narrow by construction: on most instances the
coverage row already dominates it, and it can bite only where the grouping structure,
rather than the tool count, is what forces switches.

\paragraph{Pareto-optimal cut lifting.}
When the loading dual has several optimal solutions, which a degenerate
transportation-type problem routinely does, the cuts they yield are not equally
strong. Choosing the deepest is the Pareto-optimal cut of \citet{MagnantiWong1981}:
among the alternative dual optima, select the one maximising the cut's value at a fixed
interior core point of the master. \citet{Papadakos2008} makes this practical by
observing that the core point need not be tied to the current relaxation and may simply
be updated as a convex combination of the points visited. That form is used here. A
core point is kept in the relative interior of the master polytope, initialised from
the heuristic order blended with the barycentre; at each fractional separation the
loading dual is solved a second time at the core point and its cut added alongside the
ordinary one, after which the core point is drawn halfway towards the current
relaxation point. Adding the core-point cut in addition to, never instead of, the
ordinary cut keeps the scheme valid by construction.

\subsection{Position-indexed formulations}\label{ssec:pos}

The configuration walk of Section~\ref{ssec:gtsp} is the most faithful exact picture
of the SSP, but as a model it has two flaws. There are far too many configurations to
enumerate, and forbidding the subtours that a valid walk must avoid takes exponentially
many constraints.

The position-indexed formulations, proposed by Prof.~Colares, remove both flaws by
fixing the schedule's length in advance. In place of an open-ended walk they lay out a
fixed row of \emph{positions} into which configurations are placed, and ask of each
position which configuration sits there and what it costs to change from the previous
one. Numbering the positions does two useful things at once. It makes the transition
cost a purely local quantity between adjacent slots, so the global subtour constraints
disappear, and it caps the model at a number of rows polynomial in $n$ and
$\lvert T\rvert$. That last property is the point of the construction: the row set is
fixed in advance, so column generation only ever adds columns to rows that already
exist.

\paragraph{Why not Miller--Tucker--Zemlin.}
The obvious alternative is to keep the configuration formulation and replace its
subtour constraints by Miller--Tucker--Zemlin potentials
\citep{Miller1960,DesrochersLaporte1991}. That trades one obstacle for another.

\begin{proposition}\label{prop:mtz}
A potential per configuration is exact but leaves one potential for each of
exponentially many configurations. Aggregating the potentials to one per job is
polynomial but no longer exact: it can exclude the optimum.
\end{proposition}

\begin{proof}
The aggregated form fails under job domination. If $T_i\subseteq T_j$ then every
configuration serving $j$ also serves $i$, so a transition between two such
configurations places both $i$ and $j$ at each of its ends, and the aggregated
potentials of $i$ and $j$ are forced into a two-cycle that no ordering satisfies. A
concrete instance is $b=3$ with $T=(\rset{1,4},\rset{1,3,4},\rset{1},\rset{1,2,4})$,
for which the optimum is $1$ while the aggregated model values it at $2$. This was
verified by exhaustive enumeration over covering paths.
\end{proof}

Neither variant therefore yields a tractable exact model, which is what motivates
fixing the positions instead.

\paragraph{The position--configuration formulation.}
Let $y_C^{k}\in\rset{0,1}$ indicate that configuration $C$ occupies position $k$, for
$k=1,\dots,n$, and let $w_t^{k}\ge0$ count the insertion of tool $t$ at position $k$:
\[
\begin{aligned}
  \min\ & \sum_{t,k} w_t^{k}\\
  \text{(P)}\quad & \sum_{C} y_C^{k}=1, & k&=1,\dots,n;\\
  \text{(Cov)}\quad & \sum_{k}\sum_{C\supseteq T_j} y_C^{k}\ge1, & j&\in J;\\
  \text{(W)}\quad & w_t^{k}\ge\sum_{C\ni t}\bigl(y_C^{k}-y_C^{k-1}\bigr), & t&\in T,\ k\ge2;\\
  \text{(T)}\quad & \sum_{k\ge2} w_t^{k}+\sum_{C\ni t} y_C^{1}\ge1, & t&\in\U.
\end{aligned}
\]
Constraint (P) places one configuration per position, (Cov) serves every job, (W)
charges an insertion whenever a tool enters the magazine, and (T) records that every
used tool is either present at the first position or inserted later. The
configuration variables are generated as columns; the rows number
$O(n\lvert T\rvert)$ and never change. Call this formulation PCF.

\begin{proposition}\label{prop:pcf}
PCF is exact: its integer feasible points are exactly the SSP schedules, with equal
cost. Its linear relaxation \emph{without} the counting rows (T) equals $0$. With
them it equals $q$.
\end{proposition}

The proof is in Appendix~\ref{app:proofs}.

The vanishing of the plain relaxation is the fractional pathology of
\citet{Tang1988} again: a single convex blend of configurations, repeated at every
position, satisfies (P), (Cov) and (W) at zero cost. The counting rows restore the
coverage bound, matching the multicommodity model of
Proposition~\ref{prop:sspmf-lp} with one extra family of rows.

\paragraph{Symmetry reduction.}
PCF carries a large padding symmetry. A schedule using $m<n$ distinct configurations
is encoded by repeating configurations to fill all $n$ positions, and the repeats may
sit anywhere, so one schedule has many equal-cost encodings and therefore many
distinct nodes in a search tree. Removing this freedom means relaxing (P) to an
inequality and forcing the empty positions to a suffix:
\[
  \text{(P$'$)}\quad \sum_{C} y_C^{k}\le1\ \ (\forall k);
  \qquad
  \text{(G)}\quad \sum_{C} y_C^{k+1}\le\sum_{C} y_C^{k}\ \ (k=1,\dots,n-1),
\]
with (Cov), (W) and (T) unchanged. Call this PCF$'$. At an integer point the counts
are non-increasing in $k$, so the used positions form a prefix and the empty ones a
suffix, with the cost-free first configuration at position one. The contiguity rows
(G) are not decorative: without them an interior empty position would make (W) read
the re-entry as a full reload and (T) misattribute the free load.

\begin{proposition}\label{prop:pcfprime}
PCF$'$ is exact, and its linear relaxation equals that of PCF: $q$ with the counting
rows, and $0$ without. It removes symmetric integer encodings but does not strengthen
the bound.
\end{proposition}

The proof is in Appendix~\ref{app:proofs}. The separation of concerns is deliberate:
PCF$'$ attacks symmetry, and only the next formulation attacks the bound.

\paragraph{The position--transition formulation.}
To exceed the coverage bound the formulation indexes \emph{transitions} rather than
positions. Augment the configurations with an absorbing tool-less node, used to pad
schedules that change configuration fewer than $n-1$ times, and let
$z_{C,C'}^{k}\ge0$ be a column for the step from position $k$ to position $k+1$ that
leaves $C$ and enters $C'$. The formulation is
\[
\begin{aligned}
  \min\ & \sum_{k=1}^{n-1}\sum_{(C,C')} d(C,C')\,z_{C,C'}^{k}\\
  \text{(S)}\quad & \sum_{(C,C')} z_{C,C'}^{k}=1, & k&=1,\dots,n-1;\\
  \text{(F)}\quad & \sum_{(C,C'):\,t\in C'} z_{C,C'}^{k}
                    =\sum_{(D,D'):\,t\in D} z_{D,D'}^{k+1}, & t&\in T,\ k=1,\dots,n-2;\\
  \text{(Cov)}\quad & \sum_{k=1}^{n-1}\,\sum_{(C,C'):\,T_j\subseteq C} z_{C,C'}^{k}
                    +\!\!\sum_{(C,C'):\,T_j\subseteq C'} z_{C,C'}^{\,n-1}\ \ge\ 1,
                    & j&\in J.
\end{aligned}
\]
Constraint (S) selects one transition per step. The flow-consistency rows (F) require,
tool by tool, that what enters position $k+1$ at the end of step $k$ is what leaves it
at the start of step $k+1$, so the selected columns describe a single coherent walk.
Constraint (Cov) serves every job at the configuration of some position, either the
tail of a step or the head of the last one. As in PCF the rows number
$O(n\lvert T\rvert)$ and are fixed, while the transition columns are generated. Call
this PTF.

\begin{proposition}\label{prop:ptf}
PTF is exact, and its linear relaxation is at least $q$ and can be strictly greater.
\end{proposition}

The proof is in Appendix~\ref{app:proofs}. The second claim is established by a
witness, not by a general argument, which is why it says \emph{can}.

The witnesses are the three instances of Section~\ref{ssec:results-bnp} on which the
recorded root relaxation exceeds $q$. PTF is therefore the first relaxation in this
family to improve on the elementary bound, and it does so with a fixed, polynomial row
set.

The proposition says the relaxation \emph{can} exceed $q$. It says nothing about how
often, or by how much, and Section~\ref{ssec:results-bnp} finds that on the standard
benchmark the answer to both is: rarely, and by one unit. That leaves the question that
decides whether the construction is worth pursuing.

\begin{problem}\label{prob:ptf-gap}
Is there a family of instances on which the PTF relaxation exceeds the coverage bound
$q$ by an amount that grows without bound? Equivalently: is the excess of PTF over $q$
bounded by a constant?
\end{problem}

\paragraph{Column generation and robust branching.}
Both PCF$'$ and PTF have $O(n\lvert T\rvert)$ rows but exponentially many columns, so
their relaxations are solved by column generation. Branching on a single column
$y_C^{k}$ is ineffective under the position symmetry and forces the pricing problem to
carry a growing list of forbidden columns. The branching used here is instead on the
\emph{presence} variables $a_t^{k}=\sum_{C\ni t}y_C^{k}\in\rset{0,1}$, which already
appear in (W). Fixing $a_t^{k}$ to $0$ or $1$ forbids or forces tool $t$ at position
$k$, which in the pricing problem merely shifts one dual weight and leaves the
oracle's form unchanged. This is a robust branching rule, and it is why the
Ryan--Foster rule of \citet{Ryan1981} is not needed here: that rule exists for
set-partitioning masters whose pricing problems variable branching would destroy, and
this master is not of that kind.

\paragraph{Pricing PCF$'$.}
Let $\alpha_k,\gamma_k\le0$ and $\pi_j,\mu_{t,k},\lambda_t\ge0$ be the duals of
(P$'$), (G), (Cov), (W) and (T). Collecting the coefficient of $y_C^{k}$, which enters
(W) and (T) through $a_t^{k}$, gives the reduced cost
\[
  \bar c(C,k)=\beta_k-\sum_{t\in C}\rho_t^{\,k}-\!\!\sum_{j:\,T_j\subseteq C}\!\!\pi_j,
\]
with
\[
  \rho_t^{\,k}=-\mu_{t,k}\,[k\ge2]\;+\;\mu_{t,k+1}\,[k\le n-1]\;+\;\lambda_t\,[k=1,\ t\in\U],
\]
with $\beta_k$ depending only on the position. For each $k$ the most negative reduced
cost is obtained by maximising the configuration-dependent part,
\[
  \max_{\lvert C\rvert=b}\ \sum_{t\in C}\rho_t^{\,k}
  +\!\!\sum_{j:\,T_j\subseteq C}\!\!\pi_j ,
\]
which is to say: choose $b$ tools so as to collect per-tool rewards plus a bonus for
each job whose \emph{entire} requirement lands inside the chosen set. This is a
Set-Union Knapsack Problem \citep{Goldschmidt1994}, the same class as the
grouping-pricing problem of \citet{Crama1994Column}. The branching dual enters
additively in $\rho_t^{\,k}$, so branching does not change the oracle.

\paragraph{Pricing PTF.}
For PTF a column is an ordered pair $(C,C')$, and its reduced cost contains the
bilinear term $\sum_t[t\in C'\setminus C](1-\lambda_t)$ together with chaining terms on
the tools of $C$ and $C'$. The pricing problem cannot therefore separate into a single
set: it must choose two $b$-subsets at once, with $C\ne C'$. The products
$x'_t(1-x_t)$ are linearised by the McCormick inequalities and the problem is solved as
a mixed-integer program with $2\lvert T\rvert$ binary variables. This heavier pricing
is the algorithmic price of PTF's stronger bound, and the contrast between the two
formulations --- PCF$'$ pairing a weak bound with a cheap single-set oracle, PTF a
stronger bound with a costly coupled-pair oracle --- is the central trade-off of this
section.

\paragraph{Accelerations.}
Four accelerations were implemented as independent switches, mirroring those of
Section~\ref{ssec:accel}. \emph{Heuristic pricing} answers the pricing question first
with a fast greedy subset and calls the exact oracle only to certify that no improving
column remains, so exactness is untouched while the costly oracle is skipped on all but
the final round. \emph{Multiple pricing} returns several negative-reduced-cost columns
per round instead of one. \emph{Warm starting} seeds the initial column pool from a
greedy schedule rather than the trivial one-configuration-per-job pool. \emph{Dual
stabilisation} \citep{Lubbecke2005} smooths successive dual vectors to damp the
oscillation that slows column generation on degenerate masters. None of the four
alters the relaxation value or the optimum.

\paragraph{A prototype not reported here.}
A third piece of implementation, a learned policy for choosing among the Benders
master's admissible cuts, was built during the internship but never measured on an
SSP instance. It is reported in Appendix~\ref{app:rl} and kept out of this section
for that reason; nothing elsewhere in this report rests on it.
